\documentclass[10pt,a4paper]{amsart}
\usepackage{amsmath,amssymb,amsthm,graphicx}
% Retained preamble for the Stacks Project chapter divisors.tex (Divisors).
% The chapter file pinned-original-source.tex carries only \input{preamble} on its first line, so no
% preamble text lives inside the chapter itself.  The Stacks Project uses one single shared preamble for
% every chapter; the snapshot stages it verbatim in sources/stacks-chow/preamble.tex.  The block below is
% copied mechanically, line for line, from that file's lines 14-18 (the Xy-pic setup) and lines 154-217
% (the native theorem declarations together with the mathematical macros).  Nothing was added, removed or
% redefined.  The wrapper class and the amsmath/amssymb/amsthm/graphicx base packages are supplied by the
% assembler, so the class/\externaldocument part of the shared preamble is not repeated here.
% For commutative diagrams we use Xy-pic
\usepackage[all]{xy}

% We use 2cell for 2-commutative diagrams.
\xyoption{2cell}
% Theorem environments.
%
\theoremstyle{plain}
\newtheorem{theorem}[subsection]{Theorem}
\newtheorem{proposition}[subsection]{Proposition}
\newtheorem{lemma}[subsection]{Lemma}

\theoremstyle{definition}
\newtheorem{definition}[subsection]{Definition}
\newtheorem{example}[subsection]{Example}
\newtheorem{exercise}[subsection]{Exercise}
\newtheorem{situation}[subsection]{Situation}

\theoremstyle{remark}
\newtheorem{remark}[subsection]{Remark}
\newtheorem{remarks}[subsection]{Remarks}

\numberwithin{equation}{subsection}

% Macros
%
\def\lim{\mathop{\mathrm{lim}}\nolimits}
\def\colim{\mathop{\mathrm{colim}}\nolimits}
\def\Spec{\mathop{\mathrm{Spec}}}
\def\Hom{\mathop{\mathrm{Hom}}\nolimits}
\def\Ext{\mathop{\mathrm{Ext}}\nolimits}
\def\SheafHom{\mathop{\mathcal{H}\!\mathit{om}}\nolimits}
\def\SheafExt{\mathop{\mathcal{E}\!\mathit{xt}}\nolimits}
\def\Sch{\mathit{Sch}}
\def\Mor{\mathop{\mathrm{Mor}}\nolimits}
\def\Ob{\mathop{\mathrm{Ob}}\nolimits}
\def\Sh{\mathop{\mathit{Sh}}\nolimits}
\def\NL{\mathop{N\!L}\nolimits}
\def\CH{\mathop{\mathrm{CH}}\nolimits}
\def\proetale{{pro\text{-}\acute{e}tale}}
\def\etale{{\acute{e}tale}}
\def\QCoh{\mathit{QCoh}}
\def\Ker{\mathop{\mathrm{Ker}}}
\def\Im{\mathop{\mathrm{Im}}}
\def\Coker{\mathop{\mathrm{Coker}}}
\def\Coim{\mathop{\mathrm{Coim}}}

% Boxtimes
%
\DeclareMathSymbol{\boxtimes}{\mathbin}{AMSa}{"02}

%
% Macros for moduli stacks/spaces
%
\def\QCohstack{\mathcal{QC}\!\mathit{oh}}
\def\Cohstack{\mathcal{C}\!\mathit{oh}}
\def\Spacesstack{\mathcal{S}\!\mathit{paces}}
\def\Quotfunctor{\mathrm{Quot}}
\def\Hilbfunctor{\mathrm{Hilb}}
\def\Curvesstack{\mathcal{C}\!\mathit{urves}}
\def\Polarizedstack{\mathcal{P}\!\mathit{olarized}}
\def\Complexesstack{\mathcal{C}\!\mathit{omplexes}}
% \Pic is the operator that assigns to X its picard group, usage \Pic(X)
% \Picardstack_{X/B} denotes the Picard stack of X over B
% \Picardfunctor_{X/B} denotes the Picard functor of X over B
\def\Pic{\mathop{\mathrm{Pic}}\nolimits}
\def\Picardstack{\mathcal{P}\!\mathit{ic}}
\def\Picardfunctor{\mathrm{Pic}}
\def\Deformationcategory{\mathcal{D}\!\mathit{ef}}

% --- appended 2026-10-12 (Chief): the project own definitions of the reference/slogan/history
% display environments, copied mechanically from the shared Stacks preamble lines 8-12
% (\usepackage{verbatim} plus the three \newenvironment lines, which the chapter preamble keeps
% verbatim from the project preamble). Nothing added, removed or redefined: these environments
% carry no mathematical content and the project itself typesets them as comments. Without these
% three lines a counted statement whose own text uses them fails with Environment slogan undefined.
\usepackage{verbatim}
\newenvironment{reference}{\comment}{\endcomment}
\newenvironment{slogan}{\comment}{\endcomment}
\newenvironment{history}{\comment}{\endcomment}
\usepackage{cleveref}
\crefname{originallemma}{original lemma}{original lemmas}
\Crefname{originallemma}{Original lemma}{Original lemmas}
\makeatletter
\newlabel{algebra-lemma-flat-finite-presentation-limit-flat}{{02JO}{source}{}{sourcepointer-1}{}}
\newlabel{algebra-lemma-flat-finite-presentation-limit-flat@cref}{{[originallemma][0][]02JO}{[1][0][]source}{}{}{}}
\makeatother
\crefname{originaltheorem}{original theorem}{original theorems}
\Crefname{originaltheorem}{Original theorem}{Original theorems}
\makeatletter
\newlabel{algebra-theorem-openness-flatness}{{00RC}{source}{}{sourcepointer-2}{}}
\newlabel{algebra-theorem-openness-flatness@cref}{{[originaltheorem][0][]00RC}{[1][0][]source}{}{}{}}
\makeatother
\crefname{originallemma}{original lemma}{original lemmas}
\Crefname{originallemma}{Original lemma}{Original lemmas}
\makeatletter
\newlabel{lemma-flat-at-x}{{062U}{source}{}{sourcepointer-3}{}}
\newlabel{lemma-flat-at-x@cref}{{[originallemma][0][]062U}{[1][0][]source}{}{}{}}
\makeatother
\crefname{originallemma}{original lemma}{original lemmas}
\Crefname{originallemma}{Original lemma}{Original lemmas}
\makeatletter
\newlabel{more-algebra-lemma-base-change-H1-regular}{{063Q}{source}{}{sourcepointer-4}{}}
\newlabel{more-algebra-lemma-base-change-H1-regular@cref}{{[originallemma][0][]063Q}{[1][0][]source}{}{}{}}
\makeatother
\crefname{originallemma}{original lemma}{original lemmas}
\Crefname{originallemma}{Original lemma}{Original lemmas}
\makeatletter
\newlabel{morphisms-lemma-base-change-module-flat}{{01U8}{source}{}{sourcepointer-5}{}}
\newlabel{morphisms-lemma-base-change-module-flat@cref}{{[originallemma][0][]01U8}{[1][0][]source}{}{}{}}
\makeatother
\title{A relative effective Cartier divisor over a locally finitely presented base sits in an open subscheme that is flat over the base}
\author{The Stacks Project Authors}
\date{stacks-project:divisors.tex:local-snapshot-20261012; GNU FDL 1.2 or later}
\begin{document}
\maketitle
\noindent\textit{Editorial note.} All delivered mathematical text is the author's own, copied from the pinned Stacks Project chapter file pinned-original-source.tex; no statement or proof step was rewritten, reordered or removed.\par\medskip
\noindent\textit{Editorial note.} Following the multi-claim convention for this source, source\_url is the source-level URL https://stacks.math.columbia.edu/ and the source label lemma-flat-relative-Cartier-divisor is carried in claim\_key and here rather than in a per-claim tag URL.\par\medskip
\noindent\textit{Editorial note.} The preamble is the chapter's shared Stacks Project preamble, retained as sources/stacks-divisors/retained-preamble.tex: a line-for-line excerpt of the shared Stacks preamble staged with this snapshot. No macro was invented or redefined, and this card adds no preamble text of its own.\par\medskip
\noindent\textit{Editorial note.} Statement span 4081-4088 and proof span 4090-4134 come from the halfA pool card operations/pool-cards/stacks-divisors-halfA.json and were re-read in the pinned chapter file before this card was written.\par\medskip
\noindent\textit{Editorial note.} The proof names its approximating rings with the author's own indexing and writes the two identities exhibiting the model by base change; both are preserved.\par\medskip
\noindent\textit{Editorial note.} 5 cross-references occur inside the counted spans; they are listed in external\_references as navigation pointers only.\par\medskip
\noindent\textit{Allowed standard prerequisites (editorial).}\begin{itemize}
\item The approximation of a flat and finitely presented algebra by a flat algebra of finite type over a finitely generated ring, cited from the algebra chapter. Cited at proof lines 4098-4101. Produces the Noetherian model over which the quotient by the local equation stays flat.
\item Surjectivity of the base change map on the kernel of multiplication by an element, cited from the more-algebra chapter. Cited at proof lines 4105-4107. Shows that the finitely generated kernel of multiplication by the equation over the model is killed after enlarging the base ring, since the kernel over the original ring vanishes.
\item Flatness descends along the base change defining the model, cited from the morphisms chapter. Cited at proof lines 4116-4119. Transfers the flatness statement obtained for the model back to the original morphism.
\item A relative effective Cartier divisor is flat at every point of its support, and the flat locus of a finite type morphism of Noetherian schemes is open, cited from the same chapter and the algebra chapter. Cited at proof lines 4123-4128. Gives pointwise flatness in the reduced situation and turns it into the required open neighbourhood.
\end{itemize}\medskip
\ifdefined\hypertarget\hypertarget{sourcepointer-1}{}\fi
\noindent\textit{Source reference.} Original lemma 02JO: focused-seeds/AGP/sources/stacks-snapshot-20261011/tags.txt, source lines 3343--3343. Same-chapter reference used in the retained author text at source line(s) 4106, 4234, 5430. It is supplied only as a navigation pointer so that the retained text compiles; the counted item delivers the author's own statement and the author's complete proof, and no part of the referenced item is copied in or substituted for the counted argument.\par\medskip
\ifdefined\hypertarget\hypertarget{sourcepointer-2}{}\fi
\noindent\textit{Source reference.} Original theorem 00RC: focused-seeds/AGP/sources/stacks-snapshot-20261011/tags.txt, source lines 998--998. Same-chapter reference used in the retained author text at source line(s) 4132, 4252, 5481. It is supplied only as a navigation pointer so that the retained text compiles; the counted item delivers the author's own statement and the author's complete proof, and no part of the referenced item is copied in or substituted for the counted argument.\par\medskip
\ifdefined\hypertarget\hypertarget{sourcepointer-3}{}\fi
\noindent\textit{Source reference.} Original lemma 062U: focused-seeds/AGP/sources/stacks-snapshot-20261011/tags.txt, source lines 7808--7808. Same-chapter reference used in the retained author text at source line(s) 4130, 5324. It is supplied only as a navigation pointer so that the retained text compiles; the counted item delivers the author's own statement and the author's complete proof, and no part of the referenced item is copied in or substituted for the counted argument.\par\medskip
\ifdefined\hypertarget\hypertarget{sourcepointer-4}{}\fi
\noindent\textit{Source reference.} Original lemma 063Q: focused-seeds/AGP/sources/stacks-snapshot-20261011/tags.txt, source lines 7839--7839. Same-chapter reference used in the retained author text at source line(s) 4113, 5441. It is supplied only as a navigation pointer so that the retained text compiles; the counted item delivers the author's own statement and the author's complete proof, and no part of the referenced item is copied in or substituted for the counted argument.\par\medskip
\ifdefined\hypertarget\hypertarget{sourcepointer-5}{}\fi
\noindent\textit{Source reference.} Original lemma 01U8: focused-seeds/AGP/sources/stacks-snapshot-20261011/tags.txt, source lines 2401--2401. Same-chapter reference used in the retained author text at source line(s) 4126, 5474. It is supplied only as a navigation pointer so that the retained text compiles; the counted item delivers the author's own statement and the author's complete proof, and no part of the referenced item is copied in or substituted for the counted argument.\par\medskip
\begin{lemma}
\label{lemma-flat-relative-Cartier-divisor}
Let $f : X \to S$ be a morphism of schemes.
Let $D \subset X$ be a relative effective Cartier divisor.
If $f$ is locally of finite presentation, then there exists
an open subscheme $U \subset X$ such that $D \subset U$ and
such that $f|_U : U \to S$ is flat.
\end{lemma}

\begin{proof}
Pick $x \in D$. It suffices to find an open neighbourhood $U \subset X$
of $x$ such that $f|_U$ is flat. Hence the lemma reduces to the case
that $X = \Spec(B)$ and $S = \Spec(A)$ are affine
and that $D$ is given by a nonzerodivisor $h \in B$. By assumption
$B$ is a finitely presented $A$-algebra and $B/hB$ is a flat
$A$-algebra. We are going to use absolute Noetherian approximation.

\medskip\noindent
Write $B = A[x_1, \ldots, x_n]/(g_1, \ldots, g_m)$. Assume
$h$ is the image of $h' \in A[x_1, \ldots, x_n]$. Choose a finite type
$\mathbf{Z}$-subalgebra $A_0 \subset A$ such that all the coefficients
of the polynomials $h', g_1, \ldots, g_m$ are in $A_0$. Then we can set
$B_0 = A_0[x_1, \ldots, x_n]/(g_1, \ldots, g_m)$ and $h_0$ the image
of $h'$ in $B_0$. Then $B = B_0 \otimes_{A_0} A$ and
$B/hB = B_0/h_0B_0 \otimes_{A_0} A$. By Algebra, Lemma
\ref{algebra-lemma-flat-finite-presentation-limit-flat}
we may, after enlarging $A_0$, assume that $B_0/h_0B_0$ is flat
over $A_0$. Let $K_0 = \Ker(h_0 : B_0 \to B_0)$.
As $B_0$ is of finite type over $\mathbf{Z}$ we see that $K_0$ is
a finitely generated ideal. Let $A_1 \subset A$ be a finite type
$\mathbf{Z}$-subalgebra containing $A_0$ and denote $B_1$, $h_1$, $K_1$
the corresponding objects over $A_1$. By
More on Algebra, Lemma \ref{more-algebra-lemma-base-change-H1-regular}
the map $K_0 \otimes_{A_0} A_1 \to K_1$ is surjective. On the other hand,
the kernel of $h : B \to B$ is zero by assumption. Hence every element
of $K_0$ maps to zero in $K_1$ for sufficiently large subrings
$A_1 \subset A$. Since $K_0$ is finitely generated, we conclude that
$K_1 = 0$ for a suitable choice of $A_1$.

\medskip\noindent
Set $f_1 : X_1 \to S_1$ equal to $\Spec$ of the
ring map $A_1 \to B_1$. Set $D_1 = \Spec(B_1/h_1B_1)$.
Since $B = B_1 \otimes_{A_1} A$, i.e., $X = X_1 \times_{S_1} S$,
it now suffices to prove the lemma for $X_1 \to S_1$ and the relative
effective Cartier divisor $D_1$, see
Morphisms, Lemma \ref{morphisms-lemma-base-change-module-flat}.
Hence we have reduced to the case where $A$ is a Noetherian ring.
In this case we know that the ring map $A \to B$ is flat at every
prime $\mathfrak q$ of $V(h)$ by
Lemma \ref{lemma-flat-at-x}.
Combined with the fact that the flat locus is open in this case, see
Algebra, Theorem \ref{algebra-theorem-openness-flatness}
we win.
\end{proof}

\end{document}
